\documentclass[10pt]{article}

\usepackage[margin=0.58in]{geometry}
\usepackage{array}
\usepackage{enumitem}
\usepackage[hidelinks]{hyperref}
\usepackage{xcolor}

\setlength{\parindent}{0pt}
\setlength{\parskip}{3pt}
\setlist[itemize]{leftmargin=1.1em, itemsep=1pt, topsep=2pt, parsep=0pt}
\renewcommand{\arraystretch}{1.12}

\newcommand{\projecttitle}{GeoVLM-SceneReasoner}
\newcommand{\subtitle}{Geometry-Aware Visual Reasoning for Vision-Language Models}

\begin{document}

\begin{center}
{\Large \textbf{\projecttitle}}\\[-1pt]
{\normalsize \subtitle}\\
{\small Huang Jiangzixuan \quad | \quad University of Hong Kong \quad | \quad \href{https://github.com/Hjjj0918/GeoVLM-SceneReasoner}{github.com/Hjjj0918/GeoVLM-SceneReasoner}}
\end{center}

\vspace{-5pt}

\textbf{Project Goal.}
This project studies whether vision-language models (VLMs) can reliably answer real-image spatial reasoning questions, and whether explicit object-level geometry can improve their performance. The focus is not on training a new large model, but on building a lightweight, reproducible evaluation and enhancement pipeline that separates \textit{perception failures} from \textit{reasoning failures}.

\textbf{Core Question.}
Can detection, segmentation, and monocular depth provide useful geometry signals for reasoning about object distance, relative position, support relations, occlusion, and physical size?

\vspace{2pt}
\textbf{Pipeline.}
\begin{center}
\small
\begin{tabular}{>{\centering\arraybackslash}p{0.12\linewidth}>{\centering\arraybackslash}p{0.13\linewidth}>{\centering\arraybackslash}p{0.13\linewidth}>{\centering\arraybackslash}p{0.13\linewidth}>{\centering\arraybackslash}p{0.16\linewidth}>{\centering\arraybackslash}p{0.16\linewidth}}
\textbf{Image} & \textbf{YOLO Detection} & \textbf{SAM2 Mask} & \textbf{Depth Anything V2} & \textbf{Object Geometry} & \textbf{VLM / LLM Reasoning} \\
\end{tabular}
\end{center}

\vspace{-3pt}
\begin{itemize}
  \item \textbf{Detection:} localize common scene objects and save per-image structured detections.
  \item \textbf{Segmentation:} use detection boxes as SAM2 prompts to obtain object masks.
  \item \textbf{Depth:} estimate monocular relative depth maps using Depth Anything V2.
  \item \textbf{Geometry:} summarize each object with bbox center, mask centroid, mask area, relative depth statistics, coarse position tags, and pairwise spatial relations.
  \item \textbf{Reasoning:} generate three comparable inputs: image-only VLM prompts, geometry-only LLM prompts, and GeoVLM prompts combining image and geometry.
\end{itemize}

\textbf{Current Implementation.}
The current repository implements the full preprocessing and geometry-baseline path:
\begin{itemize}
  \item real-image dataset layout, question schema, validation, and multi-view image renaming;
  \item YOLO detection, label normalization, detection visualization, SAM2 segmentation, mask visualization, and depth estimation;
  \item object-level geometry extraction and reasoning prompt generation;
  \item transparent geometry-only rule baseline and failure analysis;
  \item evaluation split generation that separates missing-target pipeline failures from geometry-available candidates.
\end{itemize}

\textbf{Preliminary Local Benchmark Status.}
The current local multi-view test set contains 37 mobile-captured images and 111 question records. The automatic split reports 81 geometry-available candidate questions and 30 upstream pipeline-failure cases caused by missing target objects. Candidate questions still require manual mask/depth review before being treated as a clean evaluation subset.

\textbf{Evaluation Design.}
\begin{center}
\small
\begin{tabular}{p{0.24\linewidth}p{0.42\linewidth}p{0.25\linewidth}}
\textbf{Track} & \textbf{Input} & \textbf{Purpose} \\
\hline
Pure VLM & image + question & image-only spatial reasoning \\
Geometry-only LLM & object-level geometry + question & value of explicit geometry alone \\
GeoVLM & image + object-level geometry + question & whether geometry improves VLM reasoning \\
\end{tabular}
\end{center}

\textbf{Future Plan.}
\begin{itemize}
  \item Build a manually reviewed clean subset by checking object masks, relative depth ordering, and question validity.
  \item Add VLM inference runners for pure VLM and GeoVLM settings, then compare against the geometry-only baseline.
  \item Expand question types to include left/right, front/back, occlusion, support, and physical-size reasoning.
  \item Evaluate on both self-captured images and public visual reasoning datasets to reduce bias from local image quality.
  \item Produce final comparison tables and failure-case analysis showing when geometry helps, when perception fails, and when reasoning remains difficult.
\end{itemize}

\vfill
{\footnotesize \textbf{Positioning:} A compact, non-robotics extension of geometry-aware multimodal reasoning, suitable for research on VLM visual reasoning, grounded perception, and lightweight benchmark construction.}

\end{document}
